\section{相关工作}
\label{sec:related_works}

\PHB{RL 后训练系统。}
大量系统工作都在解决 RL 后训练的系统瓶颈。早期框架~\cite{Harper_NeMo_a_toolkit,deepspeedchat,hu2024openrlhf} 采用静态、按阶段划分的 GPU 分区，当不同阶段负载动态变化时容易造成整体利用率下降。veRL~\cite{verl} 则通过混合控制器把多个阶段共置于同一批 GPU 上，以提高硬件效率。随后的一系列系统~\cite{liu2025specrlacceleratingonpolicyreinforcement,chen2025respecoptimizingspeculativedecoding,cheng2025fastllmposttrainingdecoupled,shao2025beatlongtaildistributionaware,RhymeRL,zhong2024rlhfuseefficientrlhftraining,areal,asyncflow,Laminar,StreamRL} 又从投机解码、阶段融合、异步训练、解耦训练等方向提出加速手段。\SystemName{}延续了解耦和异步的总体路线，但进一步引入以硬件亲和性为中心的细粒度任务映射。

\PHM{资源解耦。}
资源解耦最早可追溯到面向内存系统的体系结构工作~\cite{disaggregated-mem-isca09}，如今已经成为现代系统设计的重要基石。LegoOS~\cite{shan2018legoos}、Mira~\cite{guo2023mira} 等系统把硬件拆分成可独立管理的资源池，以提升整体利用率。类似思想在 LLM serving 中尤其常见，例如将 prefill 与 decoding 解耦~\cite{distserve, patel2023splitwise, hu2024inference, strati2024dejavu, qin2024mooncake}，甚至把注意力与 FFN 等子模块进一步拆分~\cite{megascale-infer,step3}。在 RL 后训练场景中，也有工作把 training 与 rollout 放到不同 GPU 类型上~\cite{StreamRL,seamlessflow}。相比这些系统，\SystemName{}提供了一个面向多任务 agentic RL 全生命周期的、更通用也更细粒度的解耦模型。
